% Presentación ante el CEESP, 11 de noviembre de 2026.
% Base: presentación del 22 de septiembre de 2026 (reconstruida a partir de su PDF).
% Archivo único: compila directo con pdflatex (dos pasadas); solo necesita la carpeta datos/.
% Clave de origen en cada lámina: B = base (número original), E = EXATEC, N = nueva.
\documentclass[aspectratio=169,10pt]{beamer}
% Preámbulo: reconstrucción del estilo de la presentación del 22 de septiembre de 2026.
\usepackage[T1]{fontenc}
\usepackage[utf8]{inputenc}
\usepackage[scaled=0.92]{helvet}
\renewcommand{\familydefault}{\sfdefault}
\usepackage[spanish,es-nodecimaldot,es-noshorthands]{babel}
\usepackage{graphicx,booktabs,array,tabularx,amssymb}
\usepackage{tikz,pgfplots,pgfplotstable}
\pgfplotsset{compat=1.18}
\usetikzlibrary{positioning,arrows.meta,calc,patterns}

\definecolor{marino}{HTML}{18365C}
\definecolor{oxido}{HTML}{B0541C}
\definecolor{banda}{HTML}{EEF1F6}
\definecolor{barra}{HTML}{AEB9C6}
\definecolor{medio}{HTML}{5D728D}
\definecolor{pie}{HTML}{787878}
\definecolor{rejilla}{HTML}{E5E5E5}

\usetheme{default}
\setbeamersize{text margin left=1cm,text margin right=1cm}
\setbeamertemplate{navigation symbols}{}
\setbeamercolor{normal text}{fg=black}
\setbeamercolor{structure}{fg=marino}
\setbeamercolor{frametitle}{fg=marino}
\setbeamerfont{frametitle}{series=\bfseries,size=\large}
\setbeamertemplate{frametitle}{\vspace*{0.9em}\insertframetitle\par\vspace{0.75em}{\color{black}\hrule height 0.4pt}\vspace{0.2em}}
\setbeamertemplate{footline}{\hfill{\color{pie}\scriptsize\insertframenumber\,/\,\inserttotalframenumber}\hspace*{0.45cm}\vspace*{0.35cm}}
\setbeamertemplate{itemize item}{\color{marino}--}
\setbeamertemplate{itemize subitem}{\color{marino}--}
\setbeamertemplate{enumerate item}{\color{marino}\insertenumlabel.}
\setlength{\leftmargini}{1.3em}
\setlength{\leftmarginii}{1.3em}

\newcommand{\dest}[1]{{\color{oxido}\bfseries #1}}
\newcommand{\clave}[1]{{\color{marino}\bfseries #1}}
\newcommand{\fuente}[1]{\par\vspace{0.9em}{\color{pie}\scriptsize\raggedright #1\par}}
\newcommand{\notagraf}[1]{{\color{pie}\tiny\raggedright #1\par}}
\setbeamercolor{bandacolor}{bg=banda,fg=black}
\newcommand{\banda}[1]{%
  \par\vspace{0.7em}\noindent\makebox[\textwidth][c]{\begin{beamercolorbox}[wd=0.90\textwidth,sep=5pt,center]{bandacolor}\bfseries #1\end{beamercolorbox}}\par\vspace{0.7em}}
\setbeamercolor{cajacolor}{bg=banda,fg=black}
\newcommand{\caja}[2]{%
  \begin{beamercolorbox}[wd=\linewidth,sep=5pt]{cajacolor}
    {\color{marino}\bfseries #1}\par\vspace{0.25em}\footnotesize #2
  \end{beamercolorbox}}

\pgfplotsset{
  base/.style={axis lines=left,axis line style={black},tick style={black},
    tick label style={font=\tiny},label style={font=\tiny},
    ymajorgrids,grid style={rejilla},
    /pgf/number format/.cd,fixed,1000 sep={},/tikz/.cd,
    legend style={font=\tiny,draw=none,fill=none},legend cell align=left},
}

% Metadatos directos desde LaTeX: solo título y autor.
\hypersetup{pdftitle={¿Cómo financiar la salud universal? Espacio fiscal, aseguramiento y una agenda regional},pdfauthor={Héctor Juan Villarreal Páez},pdfcreator={},pdfproducer={},pdfsubject={},pdfkeywords={}}
\pdfsuppressptexinfo=-1
\pdfinfoomitdate=1
\newcommand{\pef}{PEF 2027 (dictamen de la Cámara de Diputados)}

\title{¿Cómo financiar la salud universal?}

\begin{document}

% ===== 1. Portada (B1, adaptada) =====
\begin{frame}[plain]
  \centering
  \vspace*{1.2em}
  {\color{marino}\bfseries\Large ¿Cómo financiar la salud universal?\par}
  \vspace{0.35em}
  {\color{marino}\bfseries\normalsize Espacio fiscal, aseguramiento y una agenda regional\par}
  \vspace{2.2em}
  {\normalsize Héctor Juan Villarreal Páez\par}
  \vspace{1.1em}
  {\scriptsize Escuela de Gobierno y Transformación Pública · Tecnológico de Monterrey\\[0.3em]
  Comunidad Latinoamericana del Conocimiento sobre Financiamiento de la Salud\par}
  \vspace{1.6em}
  {\normalsize Centro de Estudios Económicos del Sector Privado (CEESP) · Sesión virtual\\[0.2em]
  11 de noviembre de 2026\par}
\end{frame}

% ===== 2. Tres preguntas (B2, adaptada) =====
\begin{frame}{Tres preguntas}
  \vspace{1.2em}
  \begin{itemize}\setlength{\itemsep}{0.35em}
    \item Me voy a concentrar en el \clave{gasto en salud} y en cómo se paga.
    \item Para financiar la salud universal hay que contestar tres preguntas: quién paga, cuánto y por cuánto tiempo.
    \item La tercera es la que menos se discute, y de ella depende que una reforma dure.
  \end{itemize}
  \vspace{0.6em}
  Voy a tocar cinco temas: lo que ya definió la demografía, el espacio fiscal, la formalización, nuestra propuesta y el financiamiento privado.
\end{frame}

% ===== 3. Menores de 15 y mayores de 65 (B3, con proyección propia) =====
\begin{frame}{México: menores de 15 y mayores de 65}
  \vspace{0.6em}
  \begin{columns}[c,onlytextwidth]
    \begin{column}{0.44\textwidth}\small
      \begin{itemize}\setlength{\itemsep}{0.7em}
        \item La población de 65 años y más pasará de \dest{13.2 millones} en 2025 a \dest{29.3 millones} en 2050.
        \item Alrededor de \dest{2036} habrá más personas de 65 y más que menores de 15 años. Naciones Unidas lo espera hasta 2051.
        \item Las personas de 80 y más se triplican: de \dest{2.7} a \dest{9.1 millones}. Es el grupo que más gasta en salud y en cuidados.
      \end{itemize}
    \end{column}
    \begin{column}{0.50\textwidth}\centering
      \begin{tikzpicture}
        \begin{axis}[base,width=\linewidth,height=4.6cm,xmin=2025,xmax=2051,ymin=0,ymax=38,
            xtick={2025,2030,2036,2040,2050},ytick={0,10,20,30},ylabel={millones de personas},
            legend style={at={(0.5,1.02)},anchor=south,legend columns=2}]
          \draw[dashed,pie] (axis cs:2036,0) -- (axis cs:2036,38);
          \addplot[barra,semithick,forget plot,restrict x to domain=2025:2050] table[col sep=comma,x=anio,y=p014]{datos/l5_ventana.csv};
          \addplot[barra,semithick,forget plot,restrict x to domain=2025:2050] table[col sep=comma,x=anio,y=p65]{datos/l5_ventana.csv};
          \addplot[marino,thick] table[col sep=comma,x=anio,y=p014]{datos/l3_nacional.csv};
          \addlegendentry{0 a 14 años}
          \addplot[oxido,thick] table[col sep=comma,x=anio,y=p65]{datos/l3_nacional.csv};
          \addlegendentry{65 y más}
          \node[font=\tiny,text=pie,anchor=south west] at (axis cs:2038,28.6) {Naciones Unidas};
        \end{axis}
      \end{tikzpicture}
    \end{column}
  \end{columns}
  \fuente{Proyección propia con datos del INEGI (Encuesta Intercensal 2025, Censo 2020, nacimientos y defunciones registrados). En gris, como referencia: Naciones Unidas, \emph{World Population Prospects 2024}, variante media.}
\end{frame}

% ===== 4. ¿Y si Naciones Unidas se queda corta? (B4, con proyección propia) =====
\begin{frame}{¿Y si Naciones Unidas se queda corta?}
  \vspace{0.3em}
  \begin{columns}[c,onlytextwidth]
    \begin{column}{0.53\textwidth}\small
      \begin{itemize}\setlength{\itemsep}{0.5em}
        \item Para 2023, Naciones Unidas estima \dest{1.91} hijos por mujer en México, y supone que en ningún año del siglo baja de \dest{1.66}.
        \item El INEGI midió \dest{1.6} en 2023, y \dest{1.23} en 2024.
        \item Fernández-Villaverde ha mostrado que las proyecciones oficiales tienden a sobreestimar la fecundidad.
      \end{itemize}
      \vspace{0.5em}
      Por eso aquí usamos una proyección propia, hecha sólo con datos del INEGI. Para 2050 da \dest{129.4 millones} de habitantes; Naciones Unidas espera 148.9.
    \end{column}
    \begin{column}{0.43\textwidth}\centering
      \begin{tikzpicture}
        \begin{axis}[base,width=\linewidth,height=4.3cm,xmin=1990,xmax=2053,ymin=1.0,ymax=3.6,
            xtick={1990,2010,2024,2050},ytick={1,1.5,2,2.5,3,3.5},ylabel={hijos por mujer}]
          \addplot[barra,thick,restrict x to domain=1990:2050] table[col sep=comma,x=anio,y=tgf]{datos/l4_fecundidad_onu.csv};
          \draw[oxido,thick,dashed] (axis cs:2024,1.23) -- (axis cs:2050,1.23);
          \node[circle,fill=oxido,inner sep=1.4pt] at (axis cs:2023,1.6) {};
          \node[circle,fill=oxido,inner sep=1.4pt] at (axis cs:2024,1.23) {};
          \node[font=\tiny,text=oxido,anchor=east] at (axis cs:2022,1.62) {INEGI 2023: 1.6};
          \node[font=\tiny,text=oxido,anchor=north east] at (axis cs:2024,1.16) {INEGI 2024: 1.23};
          \node[font=\tiny,text=oxido,anchor=south] at (axis cs:2040,1.25) {proyección propia};
          \node[font=\tiny,text=pie,anchor=south] at (axis cs:2040,1.78) {Naciones Unidas};
        \end{axis}
      \end{tikzpicture}
    \end{column}
  \end{columns}
  \fuente{ONU, WPP 2024; INEGI, ENADID 2023 y Encuesta Intercensal 2025; Fernández-Villaverde (2023). Detalle de fuentes al final.}
\end{frame}

% ===== 5. La ventana demográfica se cierra (N, con proyección propia) =====
\begin{frame}{La ventana demográfica se cierra}
  \vspace{0.5em}
  \begin{columns}[c,onlytextwidth]
    \begin{column}{0.46\textwidth}\small
      \begin{itemize}\setlength{\itemsep}{0.6em}
        \item La población de 15 a 64 años es hoy \dest{68.4\%} del total y llega a su máximo, \dest{70.1\%}, hacia \dest{2034}.
        \item Ese aumento viene sobre todo de que hay menos niños: los menores de 15 años pasan de 28.2 a 20.6 millones entre 2025 y 2034.
        \item Después baja: \dest{64.7\%} en 2050, con 29 millones de personas de 65 y más. La estructura por edades, que ayudó a las finanzas públicas durante medio siglo, deja de hacerlo en la próxima década.
      \end{itemize}
    \end{column}
    \begin{column}{0.50\textwidth}\centering
      \begin{tikzpicture}
        \begin{axis}[base,width=\linewidth,height=4.6cm,xmin=1970,xmax=2053,ymin=45,ymax=75,
            xtick={1970,1990,2010,2034,2050},ytick={50,55,60,65,70},
            ylabel={\% de la población total},clip=false]
          \draw[dashed,pie] (axis cs:2034,45) -- (axis cs:2034,70.11);
          \addplot[barra,thick,restrict x to domain=1970:2050] table[col sep=comma,x=anio,y=pet]{datos/l5_ventana.csv};
          \addplot[marino,very thick] table[col sep=comma,x=anio,y=pet]{datos/l3_nacional.csv};
          \node[circle,fill=oxido,inner sep=1.4pt] at (axis cs:2034,70.11) {};
          \node[font=\tiny\bfseries,text=oxido,anchor=south] at (axis cs:2030,70.9) {máximo: 70.1\% en 2034};
          \node[font=\tiny,text=pie,anchor=north west] at (axis cs:1985,55.5) {Naciones Unidas};
          \node[font=\tiny,text=marino,anchor=north east] at (axis cs:2049,64.3) {proyección propia};
        \end{axis}
      \end{tikzpicture}
    \end{column}
  \end{columns}
  \fuente{Proyección propia con datos del INEGI (Encuesta Intercensal 2025, Censo 2020, nacimientos y defunciones registrados). En gris, como referencia: Naciones Unidas, \emph{World Population Prospects 2024}, variante media.}
\end{frame}

% ===== 6. El ciclo de vida económico (N) =====
\begin{frame}{El ciclo de vida económico}
  \vspace{0.5em}
  \begin{columns}[c,onlytextwidth]
    \begin{column}{0.46\textwidth}\small
      \begin{itemize}\setlength{\itemsep}{0.6em}
        \item De niños y de mayores consumimos más de lo que ganamos con nuestro trabajo. Esos años se financian con \clave{transferencias}, de la familia y del gobierno, y con activos.
        \item En el perfil promedio de México, el ingreso laboral apenas alcanza al consumo alrededor de los \dest{45 años}.
        \item A los \dest{70 años}, el consumo es casi tres veces el ingreso laboral. Con 29 millones de personas de 65 y más en 2050, la vejez pesará cada vez más en esa cuenta.
      \end{itemize}
    \end{column}
    \begin{column}{0.50\textwidth}\centering
      \begin{tikzpicture}
        \begin{axis}[base,width=\linewidth,height=4.6cm,xmin=0,xmax=92,ymin=0,ymax=1.35,
            xtick={0,20,40,60,80},ytick={0,0.5,1},xlabel={edad},
            legend style={at={(0.5,1.02)},anchor=south,legend columns=2}]
          \addplot[oxido,thick] table[col sep=comma,x=edad,y=C]{datos/l6_ciclo_vida_mex2014.csv};
          \addlegendentry{Consumo}
          \addplot[marino,thick] table[col sep=comma,x=edad,y=YL]{datos/l6_ciclo_vida_mex2014.csv};
          \addlegendentry{Ingreso laboral}
        \end{axis}
      \end{tikzpicture}\par
      \notagraf{\centering México, 2014. Ingreso laboral medio de 30 a 49 años = 1. El consumo incluye el consumo público.}
    \end{column}
  \end{columns}
  \fuente{Perfiles por persona de las Cuentas Nacionales de Transferencia (National Transfer Accounts, ntaccounts.org), México, 2014. Población de 65 y más en 2050: proyección propia con datos del INEGI.}
\end{frame}

% ===== 7. También cambia la forma en que nos enfermamos (B5, sin cambios) =====
\begin{frame}{También cambia la forma en que nos enfermamos}
  \vspace{1.2em}
  \begin{itemize}\setlength{\itemsep}{0.35em}
    \item Las enfermedades no transmisibles ya explican \dest{60.1\%} de los años de vida perdidos por muerte prematura o discapacidad.
    \item Diabetes, hipertensión, cáncer: son padecimientos crónicos que se atienden durante años.
    \item El IMSS calcula que su prima de equilibrio para atención médica tendría que subir \dest{34\%} para cubrir sus gastos.
  \end{itemize}
  \vspace{0.6em}
  Esta presión va a crecer con o sin reforma. Viene de la demografía y del perfil epidemiológico del país.
  \fuente{CIEP (2026), \emph{Implicaciones del Paquete Económico 2027}, con base en IHME, OMS e IMSS.}
\end{frame}

% ===== 8. El efecto puro de la demografía sobre el gasto (N) =====
\begin{frame}{El efecto puro de la demografía sobre el gasto}
  \vspace{0.5em}
  \begin{columns}[c,onlytextwidth]
    \begin{column}{0.47\textwidth}\small
      \begin{itemize}\setlength{\itemsep}{0.6em}
        \item Multiplicamos el gasto público por persona de cada edad por la población de 2025 y por la de 2050. Nada más se mueve: no es un pronóstico.
        \item Las pensiones pasan de \dest{4.0} a \dest{6.5\%} del ingreso laboral total, y la salud de \dest{5.1} a \dest{6.1\%}. La educación baja de 6.6 a 3.7\%.
        \item El perfil es de 2004, anterior a la pensión universal: en pensiones, la cuenta se queda corta.
      \end{itemize}
      \vspace{0.4em}
      Aunque nada cambie en las reglas, el gasto sube.
    \end{column}
    \begin{column}{0.49\textwidth}\centering
      \begin{tikzpicture}
        \begin{axis}[base,width=\linewidth,height=4.5cm,ybar=1.5pt,bar width=14pt,ymin=0,ymax=8,
            ytick={0,2,4,6,8},ylabel={\% del ingreso laboral total},
            symbolic x coords={Pensiones,Salud,Educación},xtick=data,enlarge x limits=0.25,
            xtick style={draw=none},
            legend style={at={(0.5,1.02)},anchor=south,legend columns=2},area legend]
          \addplot[fill=barra,draw=marino] coordinates {(Pensiones,4.0)(Salud,5.1)(Educación,6.6)};
          \addlegendentry{2025}
          \addplot[fill=oxido,draw=oxido] coordinates {(Pensiones,6.5)(Salud,6.1)(Educación,3.7)};
          \addlegendentry{2050}
          \node[font=\tiny,anchor=south,xshift=-8pt] at (axis cs:Pensiones,4.0) {4.0};
          \node[font=\tiny\bfseries,text=oxido,anchor=south,xshift=8pt] at (axis cs:Pensiones,6.5) {6.5};
          \node[font=\tiny,anchor=south,xshift=-8pt] at (axis cs:Salud,5.1) {5.1};
          \node[font=\tiny\bfseries,text=oxido,anchor=south,xshift=8pt] at (axis cs:Salud,6.1) {6.1};
          \node[font=\tiny,anchor=south,xshift=-8pt] at (axis cs:Educación,6.6) {6.6};
          \node[font=\tiny\bfseries,text=oxido,anchor=south,xshift=8pt] at (axis cs:Educación,3.7) {3.7};
        \end{axis}
      \end{tikzpicture}\par
      \notagraf{\centering Gasto público por edad fijo (perfil de México, 2004); solo cambia la población.}
    \end{column}
  \end{columns}
  \fuente{Cálculo propio: perfiles por edad de National Transfer Accounts (ntaccounts.org), México, 2004, y proyección de población propia con datos del INEGI.}
\end{frame}

% ===== 9. Transición (B6) =====
\begin{frame}[plain]
  \vfill
  \noindent\makebox[\textwidth][c]{\begin{beamercolorbox}[wd=0.72\textwidth,sep=14pt,center]{bandacolor}
    {\color{marino}\bfseries\Large Esto ya no depende de nosotros.\par\vspace{0.5em}Lo que sigue, sí.\par}
  \end{beamercolorbox}}
  \vfill
\end{frame}

% ===== 10. ¿Cuánto queda para decidir? (B7, actualizada a enero-septiembre) =====
\begin{frame}{¿Cuánto queda para decidir?}
  \vspace{0.1em}
  \begin{columns}[c,onlytextwidth]
    \begin{column}{0.47\textwidth}\small
      \begin{itemize}\setlength{\itemsep}{0.45em}
        \item Los ingresos disponibles son \dest{20.3\%} del PIB. Dejamos fuera el petróleo: lo que entra por el Fondo Mexicano del Petróleo (0.7) sale, en buena parte, como apoyo a Pemex (0.5).
        \item Les restamos lo que no se decide año con año: gasto no programable, pensiones y la operación propia de IMSS, ISSSTE y CFE.
        \item Queda \dest{4.0\%} del PIB, más deuda, para salud, educación, infraestructura, seguridad, cuidados y todo lo demás.
      \end{itemize}
      \vspace{0.3em}
      Son sólo sumas y restas, sin supuestos de comportamiento. Y 4.0\% es el máximo.
    \end{column}
    \begin{column}{0.50\textwidth}\centering
      \begin{tikzpicture}
        \begin{axis}[base,width=\linewidth,height=4.1cm,ybar,bar width=17pt,ymin=0,ymax=24,
            ytick={0,5,10,15,20},ylabel={\% del PIB},xmin=0.4,xmax=4.6,
            xtick={1,2,3,4},xtick style={draw=none},
            xticklabels={{Ingresos\\disponibles},{menos\\no programable},{menos\\pensiones},{menos operación\\de organismos}},
            x tick label style={font=\fontsize{4.6}{5.4}\selectfont,align=center}]
          \addplot[fill=barra,draw=marino,bar shift=0pt] coordinates {(1,20.3)(2,12.1)(3,6.0)};
          \addplot[fill=oxido,draw=oxido,bar shift=0pt] coordinates {(4,4.0)};
          \node[font=\tiny,anchor=south] at (axis cs:1,20.3) {20.3};
          \node[font=\tiny,anchor=south] at (axis cs:2,12.1) {12.1};
          \node[font=\tiny,anchor=south] at (axis cs:3,6.0) {6.0};
          \node[font=\tiny\bfseries,text=oxido,anchor=south] at (axis cs:4,4.0) {4.0};
        \end{axis}
      \end{tikzpicture}\par\vspace{0.2em}
      \notagraf{Restas: no programable 8.2 (costo financiero 3.8, participaciones 4.2, Adefas 0.2); pensiones 6.1; operación propia 2.1. El resultado sin redondear es 3.95; con las cifras redondeadas la resta da 3.9.}
    \end{column}
  \end{columns}
  \fuente{Villarreal, \emph{Contabilidad de la Muerte}, actualización 2026, con SHCP, Informe enero--septiembre 2026.}
\end{frame}

% ===== 11. La deuda ya cuesta tanto como el margen de decisión (E5) =====
\begin{frame}{La deuda ya cuesta tanto como el margen de decisión}
  \centering
  \vspace{0.6em}
  \begin{tikzpicture}
    \node[draw=oxido,line width=1.2pt,fill=banda,minimum width=4.6cm,minimum height=2.7cm,align=center] (b) {
      {\small Bolsa: margen de decisión}\\[0.35em]
      {\color{oxido}\bfseries\fontsize{26}{30}\selectfont 4.0\%}\\[0.25em]
      {\small del PIB}\\[0.4em]
      {\scriptsize\color{pie} Actualización 2026}};
    \node[draw=marino,line width=1.2pt,fill=white,minimum width=4.6cm,minimum height=2.7cm,align=center,right=1.6cm of b] (d) {
      {\small Intereses de la deuda}\\[0.35em]
      {\color{marino}\bfseries\fontsize{26}{30}\selectfont 4.0\%}\\[0.25em]
      {\small del PIB}\\[0.4em]
      {\scriptsize\color{pie} PEF 2027 \;(3.8 en 2026)}};
    \node[font=\bfseries\fontsize{26}{30}\selectfont,text=marino] at ($(b.east)!0.5!(d.west)$) {=};
  \end{tikzpicture}\par
  \vspace{0.9em}
  {\normalsize Por cada peso que queda para decidir, otro peso se va en intereses.\par}
  \fuente{Bolsa y costo financiero 2026: \emph{Contabilidad de la Muerte}, actualización 2026, con datos de SHCP (enero--septiembre). Costo financiero 2027: \pef; CIEP con datos de SHCP. Razones con el PIB de SHCP; con un PIB más conservador serían algo mayores.}
\end{frame}

% ===== 12. La brecha en salud (B8, actualizada) =====
\begin{frame}{La brecha en salud}
  \vspace{0.1em}
  \begin{columns}[c,onlytextwidth]
    \begin{column}{0.47\textwidth}\small
      \begin{itemize}\setlength{\itemsep}{0.45em}
        \item En 2027 el gasto público en salud sería \dest{2.9\%} del PIB. Para llegar al 6\% que recomienda la OPS faltan \dest{3.1} puntos.
        \item Una reforma fiscal clásica completa (IVA, ISR y gastos fiscales) daría alrededor de \dest{1.5} puntos.
        \item La Miscelánea 2027 aporta alrededor de \dest{0.3} puntos.
      \end{itemize}
    \end{column}
    \begin{column}{0.50\textwidth}\centering
      \begin{tikzpicture}
        \begin{axis}[base,width=0.92\linewidth,height=3.5cm,xbar,bar width=9pt,xmin=0,xmax=3.7,
            ymajorgrids=false,xmajorgrids,axis y line*=left,y axis line style={draw=none},
            ytick={1,2,3},ytick style={draw=none},yticklabels={Miscelánea 2027,Reforma clásica,Brecha a 6\%},
            ymin=0.4,ymax=3.6,xtick={0,1,2,3},xlabel={puntos del PIB}]
          \addplot[fill=oxido,draw=oxido,bar shift=0pt] coordinates {(0.3,1)(1.46,2)(3.1,3)};
          \node[font=\tiny,anchor=west] at (axis cs:0.3,1) {$\approx$\,0.3};
          \node[font=\tiny,anchor=west] at (axis cs:1.46,2) {1.46};
          \node[font=\tiny,anchor=west] at (axis cs:3.1,3) {3.1};
        \end{axis}
      \end{tikzpicture}
    \end{column}
  \end{columns}
  \vspace{-0.2em}
  \banda{Una reforma fiscal ambiciosa, por sí sola, no alcanza para financiar la salud universal.}
  \vspace{-1.1em}
  \fuente{CIEP (2026); OPS (2014); estimación propia del paquete de reforma clásica. Miscelánea: estimación previa a las exclusiones aprobadas por el Senado; el rendimiento final es algo menor. Detalle de fuentes al final.}
\end{frame}

% ===== 13. Las pensiones le quitan espacio a la salud (B9) =====
\begin{frame}{Las pensiones le quitan espacio a la salud}
  \vspace{0.3em}
  \begin{columns}[c,onlytextwidth]
    \begin{column}{0.53\textwidth}\small
      \begin{itemize}\setlength{\itemsep}{0.45em}
        \item El PEF 2027 asigna \dest{6.3\%} del PIB a pensiones y \dest{2.9\%} a salud: más de dos pesos a pensiones por cada peso a salud.
        \item Las pensiones contributivas crecen \dest{4.7\%} real, sobre todo por IMSS e ISSSTE. Es gasto comprometido desde hace años.
        \item Si no aumentan los ingresos, cada punto adicional en pensiones es un punto menos para salud.
      \end{itemize}
      \vspace{0.4em}
      Por eso la reforma de pensiones y el financiamiento de la salud hay que verlos juntos.
    \end{column}
    \begin{column}{0.40\textwidth}\centering
      \begin{tikzpicture}
        \begin{axis}[base,width=\linewidth,height=4.2cm,ybar,bar width=20pt,ymin=0,ymax=7.5,ymajorgrids=false,
            ytick={0,2,4,6},ylabel={\% del PIB},xmin=0.4,xmax=2.6,xtick={1,2},xtick style={draw=none},
            xticklabels={Pensiones,Salud}]
          \addplot[fill=oxido,draw=oxido,bar shift=0pt] coordinates {(1,6.3)};
          \addplot[fill=barra,draw=marino,bar shift=0pt] coordinates {(2,2.9)};
          \node[font=\tiny,anchor=south] at (axis cs:1,6.3) {6.3};
          \node[font=\tiny,anchor=south] at (axis cs:2,2.9) {2.9};
        \end{axis}
      \end{tikzpicture}
    \end{column}
  \end{columns}
  \fuente{\pef. CIEP (2026), \emph{Implicaciones del Paquete Económico 2027}, caps. 5 y 10: cifras del proyecto; el dictamen incluye reasignaciones menores.}
\end{frame}

% ===== 14. ¿Financiar todo con impuestos generales? (B10) =====
\begin{frame}{¿Financiar todo con impuestos generales?}
  \vspace{0.7em}
  \begin{itemize}\setlength{\itemsep}{0.35em}
    \item Un sistema financiado sólo con impuestos generales requiere una recaudación alta y estable.
    \item México recauda alrededor de \dest{18\%} del PIB, cerca de la mitad de lo que tenían los países que lograron cobertura universal por esa vía.
    \item Además, cada peso adicional que se recaude compite con las pensiones y con el costo de la deuda.
  \end{itemize}
  \banda{Con la recaudación actual, una salud universal financiada sólo con impuestos no es viable.}
  \vspace{0.2em}
  La pregunta es cómo ampliar el aseguramiento, no cómo sustituirlo.
\end{frame}

% ===== 15. Salud en el presupuesto 2027 (B11) =====
\begin{frame}{Salud en el presupuesto 2027}
  \vspace{1.0em}
  \begin{itemize}\setlength{\itemsep}{0.35em}
    \item El gasto en salud aumenta \dest{10.6\%} real y llega a \dest{1.14 billones} de pesos. Si se descuenta la inflación de los servicios médicos, el aumento es de \dest{8.6\%}.
    \item El \dest{75.6\%} del aumento se va a población con seguridad social: IMSS crece \dest{11.6\%}, ISSSTE \dest{18.8\%} e IMSS-Bienestar \dest{9.0\%}.
    \item De ese total, \dest{760 mil millones} (\dest{67\%}) son para la población con seguridad social y \dest{376 mil millones} (\dest{33\%}) para el resto. Para el resto, el aumento es de \dest{7.6\%} real.
  \end{itemize}
  \vspace{0.6em}
  El aumento recupera parte de los recortes recientes, pero la estructura sigue igual.
  \fuente{\pef. CIEP (2026), \emph{Implicaciones del Paquete Económico 2027}, cap. 5, con base en SHCP: cifras del proyecto; el dictamen incluye reasignaciones menores.}
\end{frame}

% ===== 16. Dos sistemas y el gasto de bolsillo (B12) =====
\begin{frame}{Dos sistemas y el gasto de bolsillo}
  \vspace{0.5em}
  \begin{columns}[c,onlytextwidth]
    \begin{column}{0.44\textwidth}\small
      \begin{itemize}\setlength{\itemsep}{0.6em}
        \item El gasto público por derechohabiente del IMSS es \dest{2.1} veces el de una persona atendida por IMSS-Bienestar.
        \item El gasto de bolsillo por persona es \dest{29.3\%} mayor que el gasto público por persona en IMSS-Bienestar.
        \item Los hogares del primer decil gastan en salud \dest{3.9\%} de su ingreso; los del décimo, \dest{2.1\%}.
      \end{itemize}
    \end{column}
    \begin{column}{0.52\textwidth}\centering
      \begin{tikzpicture}
        \begin{axis}[base,width=0.9\linewidth,height=4.3cm,xbar,bar width=10pt,xmin=0,xmax=17500,
            ymajorgrids=false,axis y line*=left,y axis line style={draw=none},scaled x ticks=false,
            ytick={1,2,3,4},ytick style={draw=none},yticklabels={IMSS-Bienestar,Bolsillo,IMSS,ISSSTE},
            ymin=0.4,ymax=4.6,xtick={0,5000,10000,15000},xticklabels={0,5 mil,10 mil,15 mil},
            xlabel={pesos por persona, 2027}]
          \addplot[fill=barra,draw=marino,bar shift=0pt] coordinates {(5523,1)(11519,3)(13857,4)};
          \addplot[fill=oxido,draw=oxido,bar shift=0pt] coordinates {(7140,2)};
          \node[font=\tiny,anchor=west] at (axis cs:5523,1) {5,523};
          \node[font=\tiny,anchor=west] at (axis cs:7140,2) {7,140};
          \node[font=\tiny,anchor=west] at (axis cs:11519,3) {11,519};
          \node[font=\tiny,anchor=west] at (axis cs:13857,4) {13,857};
        \end{axis}
      \end{tikzpicture}\par
      \notagraf{\centering Bolsillo: 2024, a precios de 2027.}
    \end{column}
  \end{columns}
  \fuente{CIEP (2026), cap. 5, con base en SHCP e INEGI (ENIGH 2024).}
\end{frame}

% ===== 17. El dilema Gruber (B13) =====
\begin{frame}{El dilema Gruber}
  \vspace{0.9em}
  \begin{itemize}\setlength{\itemsep}{0.35em}
    \item En un sistema mixto es muy difícil dar el mismo servicio de salud cuando unos contribuyen y otros no.
    \item Si el servicio es igual para todos, ¿para qué cotizar? Si es distinto, terminamos con dos sistemas.
    \item Un sistema de salud tiene que lograr dos cosas: que nadie se \clave{empobrezca} por enfermarse, y que \clave{convenga} estar en el sistema contributivo.
  \end{itemize}
  \vspace{0.6em}
  En México, como en el resto de la región, se ha avanzado en una o en otra, pero rara vez en las dos al mismo tiempo.
  \fuente{El planteamiento es de Jonathan Gruber (MIT).}
\end{frame}

% ===== 18. El costo de formalizarse (B14) =====
\begin{frame}{El costo de formalizarse}
  \vspace{0.5em}
  \begin{itemize}\setlength{\itemsep}{0.3em}
    \item Al formalizarse, un trabajador empieza a pagar cuotas e impuestos y, en algunos casos, deja de recibir apoyos.
    \item El Banco Mundial encuentra que, en América Latina, la mayor parte de ese costo corresponde a las \clave{contribuciones de seguridad social}.
    \item Quién paga realmente las cuotas, el trabajador o la empresa, depende de cómo se ajustan salarios y empleo. Para México no lo tenemos bien medido.
    \item Lo que sí sabemos es que no todas las cuotas pesan igual: la \clave{cuota fija} pesa proporcionalmente mucho más en los salarios bajos.
  \end{itemize}
  \vspace{0.5em}
  Si el trabajador no valora lo que recibe a cambio, la cuota se percibe como un impuesto más. El diseño puede cambiar eso.
  \fuente{Fietz et al. (2025), \emph{(In)Formalizing Jobs in LAC}, Banco Mundial. Detalle de fuentes al final.}
\end{frame}

% ===== 19. La cuota fija pesa más donde menos se gana (N) =====
\begin{frame}{La cuota fija pesa más donde menos se gana}
  \vspace{0.5em}
  \begin{columns}[c,onlytextwidth]
    \begin{column}{0.50\textwidth}\small
      \begin{itemize}\setlength{\itemsep}{0.6em}
        \item En el seguro de enfermedades y maternidad, el patrón paga por cada trabajador una cuota fija de \dest{20.40\%} de la UMA, sin importar el salario.
        \item Para un trabajador de un salario mínimo, esa cuota equivale a \dest{7.6\%} del salario. Para uno de diez salarios mínimos, a \dest{0.8\%}.
        \item La parte de la cuota que distorsiona es la que no guarda relación con lo que el trabajador recibe, y se concentra en los salarios bajos.
      \end{itemize}
    \end{column}
    \begin{column}{0.46\textwidth}\centering
      \begin{tikzpicture}
        \begin{axis}[base,width=\linewidth,height=4.4cm,ybar,bar width=18pt,ymin=0,ymax=10,
            ytick={0,2,4,6,8},ylabel={\% del salario},xmin=0.4,xmax=4.6,
            xtick={1,2,3,4},xtick style={draw=none},xticklabels={1,2,5,10},xlabel={salarios mínimos}]
          \addplot[fill=oxido,draw=oxido,bar shift=0pt] coordinates {(1,7.60)};
          \addplot[fill=barra,draw=marino,bar shift=0pt] coordinates {(2,3.80)(3,1.52)(4,0.76)};
          \node[font=\tiny\bfseries,text=oxido,anchor=south] at (axis cs:1,7.60) {7.6};
          \node[font=\tiny,anchor=south] at (axis cs:2,3.80) {3.8};
          \node[font=\tiny,anchor=south] at (axis cs:3,1.52) {1.5};
          \node[font=\tiny,anchor=south] at (axis cs:4,0.76) {0.8};
        \end{axis}
      \end{tikzpicture}\par
      \notagraf{\centering Cuota fija patronal como \% del salario. Valores de 2026: UMA de 117.31 pesos diarios y salario mínimo general de 315.04.}
    \end{column}
  \end{columns}
  \fuente{Ley del Seguro Social, art. 106, fr. I; IMSS, \emph{Informe al Ejecutivo Federal y al Congreso de la Unión 2025--2026}; INEGI, UMA 2026; CONASAMI, salario mínimo general 2026. Cálculos propios.}
\end{frame}

% ===== 20. Cómo decide un trabajador (N) =====
\begin{frame}{Cómo decide un trabajador}
  \vspace{0.1em}
  \begin{columns}[T,onlytextwidth]
    \begin{column}{0.49\textwidth}
      \caja{Lo que cuesta ser formal}{%
        Cuotas de seguridad social: en América Latina son el componente mayor.\par\vspace{0.3em}
        Impuestos sobre el ingreso.\par\vspace{0.3em}
        Apoyos y servicios no contributivos que se pueden perder.}
    \end{column}
    \begin{column}{0.49\textwidth}
      \caja{Lo que recibe a cambio}{%
        Salud, pensión, vivienda, guarderías, incapacidades.\par\vspace{0.3em}
        Cuenta lo que el trabajador cree que vale, no lo que le cuesta al sistema.\par\vspace{0.3em}
        Si la cobertura sin cotizar se parece a la contributiva, cotizar vale menos.}
    \end{column}
  \end{columns}
  \vspace{0.5em}
  {\small
  \begin{itemize}\setlength{\itemsep}{0.25em}
    \item El costo neto de formalizarse es lo que se paga menos el valor que el trabajador asigna a lo que recibe. Para un asalariado de la región, va de \dest{20} a \dest{40\%} del ingreso del hogar.
    \item De ahí salen dos reglas de diseño: lo que se recibe debe depender de lo que se cotiza, y la cuota por el mismo ingreso no debe depender de si el trabajo es formal o informal.
    \item Todavía no hemos estimado cuánto responde la formalidad en México a estos cambios. Lo que sigue es lógica de diseño, no un resultado.
  \end{itemize}}
  \vspace{-0.5em}
  \fuente{Fietz et al. (2025), \emph{(In)Formalizing Jobs in LAC}, Banco Mundial, cap. 5: mediana en seis países; en México los impuestos pesan más que en el resto.}
\end{frame}

% ===== 21. El escalón (N) =====
\begin{frame}{El escalón}
  \vspace{0.5em}
  \begin{columns}[c,onlytextwidth]
    \begin{column}{0.50\textwidth}\small
      \begin{itemize}\setlength{\itemsep}{0.6em}
        \item Al pasar a un empleo formal, el ingreso neto puede caer: se empiezan a pagar cuotas y se pueden perder apoyos. Ese salto es la \clave{trampa de la formalización}.
        \item El escalón es más alto en los ingresos medios y bajos, donde la cuota fija pesa más y los apoyos importan más.
        \item Uruguay reconoció el desincentivo. Desde 2022 suspendió el control del tope de ingreso de sus asignaciones familiares, y entraron \dest{21 mil} hogares adicionales.
      \end{itemize}
    \end{column}
    \begin{column}{0.46\textwidth}\centering
      \begin{tikzpicture}[x=0.5cm,y=0.5cm,font=\tiny]
        \draw[-{Straight Barb[length=2pt]},pie] (0,0) -- (0,6.4);
        \draw[-{Straight Barb[length=2pt]},pie] (0,0) -- (10.6,0);
        \node[anchor=south] at (0.6,6.4) {ingreso neto};
        \node[anchor=north east] at (10.6,-0.1) {ingreso formal};
        \draw[marino,thick,dashed] (0,2) -- (10,5.6);
        \draw[oxido,thick] (0,2.6) -- (4.5,4.0);
        \draw[oxido,thick] (4.5,2.9) -- (10,4.9);
        \draw[oxido,thick,densely dotted] (4.5,4.0) -- (4.5,2.9);
        \node[text=oxido,anchor=west] at (4.7,2.3) {el escalón};
        \node[text=marino,anchor=south east] at (9.8,5.7) {sin saltos};
      \end{tikzpicture}\par\vspace{0.6em}
      \notagraf{\centering Esquema conceptual, sin escala.}
    \end{column}
  \end{columns}
  \fuente{Fietz et al. (2025), Banco Mundial; BPS Uruguay, resoluciones 2021--2026; MEF Uruguay, Rendición de Cuentas 2023. Detalle de fuentes al final.}
\end{frame}

% ===== 22 (antes 23). Infonavit: la propuesta (N) =====
\begin{frame}{Infonavit: la propuesta}
  \vspace{0.1em}
  \begin{columns}[T,onlytextwidth]
    \begin{column}{0.49\textwidth}
      \caja{Qué se propone}{%
        Eliminar la aportación patronal de \textbf{5\%} del salario base de cotización al Infonavit.\par\vspace{0.3em}
        Destinar \textbf{2 puntos} a un seguro solidario de salud y cuidados en la vejez: un fondo común, no cuentas individuales. Primero cubre el déficit de los gastos médicos de pensionados; lo que sobra va a cuidados de largo plazo.\par\vspace{0.3em}
        Los otros \textbf{3 puntos} se dejarían de cobrar: la tasa legal sobre la nómina baja 3 puntos.}
    \end{column}
    \begin{column}{0.49\textwidth}
      \caja{El Instituto y los trabajadores}{%
        El Infonavit conserva su patrimonio y sus obligaciones como mutualidad.\par\vspace{0.3em}
        El trabajador deja de acumular 5\% de su salario en la subcuenta de vivienda y recibe a cambio un seguro.\par\vspace{0.3em}
        Quien no usa crédito tendría menos saldo para su pensión. Ese efecto no lo hemos cuantificado, y tampoco el efecto sobre la formalidad.}
    \end{column}
  \end{columns}
  \vspace{0.5em}
  {\small El seguro de gastos médicos de pensionados del IMSS ya es deficitario, y el envejecimiento presiona justo ahí.\par}
  \vspace{-0.4em}
  \fuente{Ejercicio propio, preliminar; no es una propuesta del IMSS ni del Gobierno Federal. Son tasas legales: no miden quién soporta la cuota. Antecedente: Villarreal (2022), «El Infonavit, ¿qué sigue?», \emph{Arena Pública}.}
\end{frame}

% ===== Bismarck, Beveridge y nuestra propuesta (B15) =====
\begin{frame}{Bismarck, Beveridge y nuestra propuesta}
  \vspace{-0.2em}
  {\scriptsize
  \begin{tabularx}{\textwidth}{@{}>{\raggedright\arraybackslash}p{2.1cm}>{\raggedright\arraybackslash}p{2.0cm}>{\raggedright\arraybackslash}p{1.6cm}>{\raggedright\arraybackslash}X@{}}
    \toprule
    \textbf{Decisión} & \textbf{Bismarck} & \textbf{Beveridge} & \textbf{Nuestra propuesta}\\
    \midrule
    ¿Con qué se financia? & Cuotas obrero-patronales & Impuestos generales & Se mantienen las cuotas obrero-patronales; quien no cotiza paga una prima según su ingreso; el gobierno completa\\\addlinespace
    ¿La prima es de cada persona? & Sí & No & Sí\\\addlinespace
    ¿Quién la paga? & Trabajador y patrón (en México, también el gobierno) & El gobierno & Asalariados, igual que hoy. Los demás, según su ingreso; el gobierno paga la prima completa sólo a los más pobres\\\addlinespace
    ¿El derecho depende del pago? & Sí & No, basta con residir & Sí, y todos tienen su prima cubierta\\
    \bottomrule
  \end{tabularx}}
  \vspace{-0.3em}
  \banda{Nuestra propuesta combina de otra manera las mismas cuatro decisiones.}
\end{frame}

% ===== Cómo funcionaría (B16, adaptada) =====
\begin{frame}{Cómo funcionaría}
  \vspace{0.5em}
  \begin{itemize}\setlength{\itemsep}{0.35em}
    \item Quien tiene empleo formal sigue pagando sus \clave{cuotas obrero-patronales}, como hoy. Lo que cambia es que debe ver con claridad qué paga y qué recibe.
    \item Quien no cotiza se inscribe en los seguros médicos del IMSS y paga una \clave{prima según su ingreso}. El gobierno completa la diferencia.
    \item Sólo a la población más pobre el gobierno le paga la prima completa, como un impuesto negativo. Si la pagara para todos, se perdería el incentivo a cotizar.
  \end{itemize}
  \vspace{-0.2em}
  \banda{Todos dentro del mismo sistema de aseguramiento, sin esquemas paralelos.}
  \vspace{-0.9em}
  \banda{\small La propuesta no sube las cuotas obrero-patronales. Con el cambio en Infonavit, la tasa legal sobre la nómina bajaría 3 puntos.}
\end{frame}

% ===== El IMSS como sistema nacional de seguros de salud (B17) =====
\begin{frame}{El IMSS como sistema nacional de seguros de salud}
  \vspace{0.9em}
  \begin{itemize}\setlength{\itemsep}{0.3em}
    \item Partimos de lo que existe. El IMSS tiene los hospitales, el personal, buena parte de la formación de especialistas y el registro de derechohabientes. Para la población sin seguridad social, en cambio, el modelo cambió tres veces entre \dest{2019} y \dest{2023}.
    \item Que el IMSS agrupe los servicios de salud ligados al empleo formal: ISSSTE, Pemex e incluso fuerzas armadas.
    \item Que se dedique sólo a salud. Pensiones y cuidados, en otras instancias.
    \item Que pueda funcionar como caja de compensación y contratar servicios privados con recursos públicos.
  \end{itemize}
  \fuente{Diario Oficial de la Federación, 29-XI-2019, 31-VIII-2022 y 29-V-2023; Villarreal (2024). Detalle de fuentes al final.}
\end{frame}

% ===== Consolidar transferencias (B18) =====
\begin{frame}{Consolidar transferencias}
  \vspace{0.6em}
  \begin{columns}[c,onlytextwidth]
    \begin{column}{0.47\textwidth}\small
      \begin{itemize}\setlength{\itemsep}{0.6em}
        \item Los programas con límites de ingreso generan un escalón: quien rebasa el límite pierde el apoyo completo.
        \item Con un \clave{impuesto negativo sobre la renta}, el apoyo baja poco a poco conforme aumenta el ingreso, y formalizarse siempre conviene.
        \item Esto requiere medir bien los ingresos. La factura electrónica, los pagos digitales y el RFC como padrón único lo hacen posible.
      \end{itemize}
    \end{column}
    \begin{column}{0.49\textwidth}\centering
      \begin{tikzpicture}[x=1cm,y=1cm,font=\tiny]
        \draw[-{Straight Barb[length=2pt]},pie] (0,0) -- (0,2.7);
        \draw[-{Straight Barb[length=2pt]},pie] (0,0) -- (4.9,0);
        \node[anchor=south] at (0,2.75) {beneficio};
        \node[anchor=west] at (4.95,0) {ingreso};
        \draw[oxido,thick] (0.2,2.15) -- (2.2,2.15);
        \draw[oxido,thick,dashed] (2.2,2.15) -- (2.2,0.3);
        \draw[oxido,thick] (2.2,0.3) -- (4.5,0.3);
        \draw[marino,thick] (0.2,1.85) -- (4.5,0.15);
        \node[text=oxido,anchor=south west] at (0.55,2.17) {sistema actual};
        \node[text=marino,anchor=south west] at (2.8,0.95) {impuesto negativo};
        \node[text=pie,anchor=north] at (2.2,-0.08) {umbral};
      \end{tikzpicture}\par\vspace{0.6em}
      \notagraf{\centering Esquema conceptual, sin escala.}
    \end{column}
  \end{columns}
\end{frame}

% ===== ¿De dónde saldría el 6 % del PIB? (B19) =====
\begin{frame}{¿De dónde saldría el 6\% del PIB?}
  \vspace{0.1em}
  \begin{columns}[c,onlytextwidth]
    \begin{column}{0.48\textwidth}\small
      \begin{itemize}\setlength{\itemsep}{0.4em}
        \item \clave{Cuotas de salud:} de \dest{1.5\%} a \dest{2.5\%} del PIB, con el crecimiento de la masa salarial y la formalización.
        \item \clave{Gobierno federal}, sin contar cuotas: de \dest{1.4\%} a \dest{2.5\%}, sumando \dest{0.1} puntos cada año. La fuente serían los gastos fiscales.
        \item \clave{Estados y Ciudad de México:} \dest{1\%} del PIB, con participaciones e ingresos propios, sobre todo predial.
      \end{itemize}
      \vspace{0.3em}
      Nada de esto funciona si no se contiene el gasto en pensiones.
    \end{column}
    \begin{column}{0.48\textwidth}\centering
      \begin{tikzpicture}
        \begin{axis}[base,width=0.95\linewidth,height=4.2cm,ybar stacked,bar width=26pt,ymin=0,ymax=7.3,
            ytick={0,2,4,6},ylabel={\% del PIB},xmin=0.4,xmax=2.6,xtick={1,2},xtick style={draw=none},
            xticklabels={Hoy,Meta 2036},
            legend style={at={(0.5,-0.2)},anchor=north,legend columns=3},area legend]
          \addplot[fill=barra,draw=marino] coordinates {(1,1.5)(2,2.5)};
          \addlegendentry{Cuotas}
          \addplot[fill=medio,draw=marino] coordinates {(1,1.4)(2,2.5)};
          \addlegendentry{Federal}
          \addplot[fill=oxido,draw=oxido] coordinates {(1,0)(2,1.0)};
          \addlegendentry{Estados}
          \node[font=\tiny\bfseries,text=oxido,anchor=south] at (axis cs:1,2.9) {2.9};
          \node[font=\tiny\bfseries,text=oxido,anchor=south] at (axis cs:2,6.0) {6.0};
        \end{axis}
      \end{tikzpicture}\par
      \notagraf{\centering Hoy: aportación estatal propia no desglosada.}
    \end{column}
  \end{columns}
  \fuente{Villarreal (2024), «¿Se puede reformar la seguridad social? Parte II», \emph{Arena Pública}.}
\end{frame}

% ===== Modalidad 40 (N; fusión de EXATEC 12 y 13) =====
\begin{frame}{Seguir cotizando es razonable; comprar una pensión no lo es}
  \vspace{0.1em}
  \begin{columns}[T,onlytextwidth]
    \begin{column}{0.49\textwidth}
      \caja{Modalidad 40: lo que cambió}{%
        \textbf{Hasta el 21 de octubre.} Cualquier salario de cotización de hasta 25 UMA (casi 3 mil pesos diarios), sin relación con el salario previo, y con inscripciones con fecha anterior.\par\vspace{0.35em}
        \textbf{Desde el 22 de octubre.} El último salario registrado, o ese salario más una UMA (unos 117 pesos diarios), nunca por debajo del mínimo. Sin inscripciones retroactivas; trámite solo digital.\par\vspace{0.35em}
        Los inscritos conservan sus condiciones hasta terminar su periodo. Es una medida temporal: seis meses para estimar el costo y otros seis para proponer cambios a las normas.}
    \end{column}
    \begin{column}{0.49\textwidth}
      \caja{Quién paga}{%
        Para quien cotizó antes de julio de 1997, la pensión usa el salario promedio de los últimos cinco años: la estrategia abría la puerta a pensiones vitalicias de más de 80 mil pesos al mes.\par\vspace{0.35em}
        Las cuotas no alcanzan; la diferencia la paga el Gobierno Federal, es decir, todos los contribuyentes, incluidos quienes nunca tendrán pensión.\par\vspace{0.35em}
        Solo está al alcance de quien puede adelantar cientos de miles de pesos.}
    \end{column}
  \end{columns}
  \vspace{0.5em}
  {\small Nada de la propuesta de salud funciona si no se contiene el gasto en pensiones. Este es un primer paso, y necesita base legal firme.\par}
  \vspace{-0.4em}
  \fuente{Acuerdo ACDO.AS2.HCT.101026/302.P.DG, \emph{Diario Oficial de la Federación}, 21 de octubre de 2026; Ley del Seguro Social de 1973.}
\end{frame}

% ===== Financiamiento privado (I) (B20, primera mitad) =====
\begin{frame}{Financiamiento privado (I): prepago, no bolsillo}
  \vspace{0.5em}
  \begin{columns}[c,onlytextwidth]
    \begin{column}{0.52\textwidth}
      \begin{itemize}\setlength{\itemsep}{0.5em}
        \item Con una meta de \dest{8\%} del PIB en salud hacia 2036, \dest{6} puntos serían públicos y \dest{2} privados: alrededor de \dest{25\%} del total, contra la mitad o más de hoy.
        \item Lo importante es que el gasto privado sea \clave{prepago}, y no gasto de bolsillo por urgencias.
        \item El prepago reparte el riesgo entre muchos y se paga antes de enfermarse. El gasto de bolsillo llega cuando el hogar menos puede.
      \end{itemize}
    \end{column}
    \begin{column}{0.42\textwidth}\centering
      \begin{tikzpicture}[x=1cm,y=0.36cm]
        \fill[barra] (0,0) rectangle (2.6,6); \draw[marino] (0,0) rectangle (2.6,6);
        \fill[oxido] (0,6) rectangle (2.6,8);
        \node[font=\small] at (1.3,3) {Público: 6};
        \node[font=\small,text=white] at (1.3,7) {Privado: 2};
        \node[font=\small\bfseries,text=oxido,anchor=south] at (1.3,8) {8\% del PIB};
        \node[font=\scriptsize,anchor=north] at (1.3,-0.1) {Meta 2036};
        \draw (-0.3,0) -- (2.9,0);
      \end{tikzpicture}
    \end{column}
  \end{columns}
  \fuente{Villarreal (2024), «¿Se puede reformar la seguridad social? Parte II», \emph{Arena Pública}.}
\end{frame}

% ===== Financiamiento privado (II) (N, de la segunda mitad de B20) =====
\begin{frame}{Financiamiento privado (II): instrumentos y reglas}
  \vspace{0.6em}
  \begin{itemize}\setlength{\itemsep}{0.4em}
    \item \clave{Consultorios adjuntos a farmacias:} regularlos y conectarlos con el resto del sistema.
    \item \clave{Seguros médicos privados de bajo costo,} con estímulos fiscales o descuentos en la prima del IMSS.
    \item \clave{Seguros de largo plazo,} que se paguen durante la vida laboral y cubran en la vejez. Sigue siendo una pregunta abierta.
  \end{itemize}
  \vspace{0.2em}
  \banda{¿Qué reglas necesita el sector privado para participar sin fragmentar el sistema?}
  \vspace{-0.9em}
  \fuente{Villarreal (2024), «¿Se puede reformar la seguridad social? Parte II», \emph{Arena Pública}.}
\end{frame}

% ===== Dos preguntas para cualquier propuesta (B21) =====
\begin{frame}{Dos preguntas para cualquier propuesta}
  \vspace{1.6em}
  \begin{enumerate}\setlength{\itemsep}{0.35em}
    \item ¿Se puede pagar durante toda la transición demográfica sin que la deuda se vuelva insostenible?
    \item ¿Fortalece o debilita el incentivo a cotizar?
  \end{enumerate}
  \vspace{0.6em}
  Si una propuesta falla en alguna de las dos, tiene un problema de fondo. Esto aplica a lo que propongo y a cualquier alternativa. Para contestar la primera hay que hacer la cuenta completa: salud, pensiones y reforma fiscal juntas.
\end{frame}

% ===== Sin ancla, cualquier reforma es frágil (E15 v2, adaptada) =====
\begin{frame}{El techo de deuda no ancla: la LIF autoriza casi un punto más de lo programado}
  \centering
  \vspace{0.2em}
  \begin{tikzpicture}
    \begin{axis}[base,width=0.74\textwidth,height=3.3cm,xbar,bar width=11pt,xmin=0,xmax=5.8,
        ymajorgrids=false,axis y line*=left,y axis line style={draw=none},
        ytick={1,2},ytick style={draw=none},
        yticklabels={{Endeudamiento programado\\(LIF 2027, art. 1o)},{Techo autorizado\\(LIF 2027, art. 2o)}},
        y tick label style={font=\scriptsize,align=right},
        ymin=0.35,ymax=2.65,xtick={0,1,2,3,4,5},xlabel={\% del PIB},clip=false]
      \addplot[fill=barra,draw=marino,bar shift=0pt] coordinates {(4.0,1)};
      \addplot[fill=oxido,draw=oxido,bar shift=0pt] coordinates {(4.9,2)};
      \node[font=\small\bfseries,anchor=west] at (axis cs:4.0,1) {4.0};
      \node[font=\small\bfseries,text=oxido,anchor=west] at (axis cs:4.9,2) {4.9};
      \draw[marino,dashed,thick] (axis cs:3.9,0.35) -- (axis cs:3.9,2.75);
      \node[font=\tiny,text=marino,anchor=south] at (axis cs:3.9,2.75) {Meta de RFSP: 3.9};
    \end{axis}
  \end{tikzpicture}\par
  \vspace{0.1em}
  {\color{pie}\scriptsize Gobierno Federal. LIF 2027, arts. 1o y 2o; SHCP, CGPE 2027. \% del PIB con el PIB de SHCP.\par}
  \vspace{0.1em}
  \banda{\small Responder la primera pregunta exige un ancla fiscal creíble: la reforma de salud y la de ingresos deben ir amarradas a ella.}
\end{frame}

% ===== La Comunidad (B22 y B23 fusionadas) =====
\begin{frame}{La Comunidad}
  \vspace{0.5em}
  {\small
  \begin{itemize}\setlength{\itemsep}{0.4em}
    \item En la región se ha discutido mucho cómo ampliar la \clave{cobertura} en salud, y mucho menos cómo \clave{pagarla} en los próximos veinte o treinta años. De ahí surge la Comunidad.
    \item La fundan tres instituciones:
      \begin{itemize}
        \item Tecnológico de Monterrey --- coordinación, Héctor Villarreal
        \item Universidad de los Andes, Colombia --- Ramiro Guerrero
        \item Pontificia Universidad Católica de Chile --- Pablo Celhay
      \end{itemize}
    \item La primera reunión de trabajo fue en \clave{San José, Costa Rica, del 5 al 7 de octubre}.
    \item Dos de las preguntas abiertas:
      \begin{itemize}
        \item ¿Qué parte del gasto de bolsillo podría convertirse en prepago?
        \item ¿Con qué reglas el financiamiento privado complementa al público en lugar de debilitarlo?
      \end{itemize}
  \end{itemize}}
\end{frame}

% ===== Cierre (B24) =====
\begin{frame}[plain]
  \vfill
  \noindent\makebox[\textwidth][c]{\begin{beamercolorbox}[wd=0.78\textwidth,sep=10pt,center]{bandacolor}
    {\color{marino}\bfseries\normalsize ¿Podemos diseñar un sistema de salud que la gente valore lo suficiente para querer participar, y que el país pueda pagar durante su transición demográfica?\par}
  \end{beamercolorbox}}
  \begin{center}
  \vspace{1.6em}
  {\small Héctor Juan Villarreal Páez · hjvp@itesm.mx\\[0.2em]
  Escuela de Gobierno y Transformación Pública · Tecnológico de Monterrey\par}
  \end{center}
  \vfill
\end{frame}

% ===== Fuentes (B25-26, actualizadas) =====
\begin{frame}{Fuentes (1/2)}
  \vspace{0.2em}
  {\fontsize{6.6}{8.4}\selectfont\raggedright
  \textbf{Demografía.} Proyección propia de la población de México, 2025--2050, por edad simple y sexo, hecha sólo con datos del INEGI: Encuesta Intercensal 2025 (población base, fecundidad de 1.23 hijos por mujer en 2024 y emigración), Censo de Población y Vivienda 2020, nacimientos y defunciones registrados; patrón de mortalidad por edad del IBGE ajustado a las defunciones registradas por el INEGI. La fecundidad se mantiene en 1.23. Referencia: Naciones Unidas, DESA, \emph{World Population Prospects 2024}, variante media. INEGI, Encuesta Nacional de la Dinámica Demográfica 2023. Fernández-Villaverde, J. (2023), \emph{The Demographic Future of Humanity: Facts and Consequences}.\par\vspace{0.6em}
  \textbf{Ciclo de vida y efecto de la demografía sobre el gasto.} National Transfer Accounts (ntaccounts.org), perfiles por edad de México, 2004 y 2014. Gasto con perfiles fijos: cálculo propio con el perfil de 2004 y una proyección de población propia con datos del INEGI; resultados en \% del ingreso laboral total, no del PIB.\par\vspace{0.6em}
  \textbf{Presupuesto 2027, salud y pensiones.} PEF 2027: dictamen de la Cámara de Diputados; gasto neto total de 10.64 billones de pesos. CIEP (2026), \emph{Implicaciones del Paquete Económico 2027}, 11 de septiembre de 2026: gasto en salud, gasto per cápita por subsistema, gasto de bolsillo, inflación de salud, carga de enfermedad (IHME, 2021; OMS, 2024), prima de equilibrio del IMSS (IMSS, 2024), pensiones, Miscelánea e ingresos petroleros; sus cifras corresponden al proyecto de presupuesto. Meta de 6\% del PIB: OPS (2014), \emph{Estrategia para el acceso universal a la salud y la cobertura universal de salud}.\par\vspace{0.6em}
  \textbf{Espacio disponible.} Villarreal, \emph{Contabilidad de la Muerte}, actualización con el Informe de Finanzas Públicas y Deuda Pública enero--septiembre 2026 (SHCP) y los Criterios Generales de Política Económica 2027. PIB nominal de 36,500 mmdp; el supuesto favorece a la cuenta oficial. CIEP (2026) estima un espacio fiscal de 1.6\% del PIB con una definición distinta.\par\vspace{0.6em}
  \textbf{Deuda y reforma al ISR.} LIF 2027, aprobada por el Senado el 27 de octubre de 2026, arts. 1o, 2o y 4o. SHCP, Criterios Generales de Política Económica 2027. Congreso de la Unión, minuta de reforma a la LISR 2027.\par}
\end{frame}

\begin{frame}{Fuentes (2/2)}
  \vspace{0.2em}
  {\fontsize{6.6}{8.4}\selectfont\raggedright
  \textbf{Reforma fiscal clásica.} Estimación propia del rendimiento de un paquete de reforma clásica (IVA, ISR y gastos fiscales), alrededor de 1.5 puntos del PIB.\par\vspace{0.6em}
  \textbf{Trampa de la formalización.} Fietz, Joubert, Ñopo, Ocampo, Packard, Posadas y Rodriguez Chamussy (2025), \emph{(In)Formalizing Jobs in Latin America and the Caribbean: Taxes, Benefits, and Labor Market Incentives}, Banco Mundial, cap. 5. Uruguay: Banco de Previsión Social, resoluciones de directorio 2021--2026; Ministerio de Economía y Finanzas, Rendición de Cuentas 2023.\par\vspace{0.6em}
  \textbf{Cuota fija.} Ley del Seguro Social, art. 106, fr. I. IMSS, \emph{Informe al Ejecutivo Federal y al Congreso de la Unión sobre la situación financiera y los riesgos del IMSS 2025--2026}: cuota fija patronal de 20.40\% de la UMA. INEGI, UMA 2026: 117.31 pesos diarios. CONASAMI, salario mínimo general 2026: 315.04 pesos diarios.\par\vspace{0.6em}
  \textbf{Infonavit.} Ejercicio propio, preliminar, sobre tasas legales. No es una propuesta del IMSS ni del Gobierno Federal. Antecedente: Villarreal, H. (2022), «El Infonavit, ¿qué sigue?», \emph{Arena Pública}, 20 de diciembre de 2022.\par\vspace{0.6em}
  \textbf{Modalidad 40.} Acuerdo ACDO.AS2.HCT.101026/302.P.DG del Consejo Técnico del IMSS, \emph{Diario Oficial de la Federación}, 21 de octubre de 2026. Ley del Seguro Social de 1973.\par\vspace{0.6em}
  \textbf{Modelo de aseguramiento y metas de financiamiento.} Villarreal, H. (2024), «¿Se puede reformar la seguridad social? Parte II», \emph{Arena Pública}, 10 de septiembre de 2024. Dilema de Gruber: Jonathan Gruber (MIT).\par\vspace{0.6em}
  \textbf{Arquitectura institucional.} Diario Oficial de la Federación: 29-XI-2019 (Insabi), 31-VIII-2022 (IMSS-Bienestar), 29-V-2023 (reforma a la Ley General de Salud).\par\vspace{0.6em}
  \textbf{Comunidad.} Documento de presentación de la iniciativa; reunión de trabajo en San José, Costa Rica, 5--7 de octubre de 2026.\par}
\end{frame}

% ===== A1. El paquete 2027 mueve la base en el margen (E6, actualizada) =====
\begin{frame}{Apéndice · El paquete 2027 mueve la base en el margen}
  \vspace{0.2em}
  \begin{columns}[T,onlytextwidth]
    \begin{column}{0.49\textwidth}
      \caja{Amplía la base (en el margen)}{%
        Deducciones de ISR topadas en 96.67\% de los ingresos acumulables, para empresas con ingresos mayores a 50 mdp: un impuesto mínimo de alrededor de 1\% de los ingresos.\par\vspace{0.35em}
        El Senado (27 de octubre) excluyó a las armadoras y ensambladoras de vehículos en depósito fiscal y a los expendedores de gasolina y diésel. La Cámara de Diputados ratificó los cambios el 4 de noviembre.\par\vspace{0.35em}
        Pérdidas fiscales: amortización hasta 50\% de la utilidad fiscal del ejercicio.}
    \end{column}
    \begin{column}{0.49\textwidth}
      \caja{Abarata el incumplimiento}{%
        Programa de regularización con condonación de 100\% de multas, recargos y gastos de ejecución.\par\vspace{0.35em}
        Ampliado a grandes contribuyentes: se eliminó el tope de 300 mdp de ingresos.}
    \end{column}
  \end{columns}
  \vspace{0.6em}
  {\small Las exclusiones reducen el rendimiento de la medida. No cambia la restricción de fondo.\par}
  \vspace{-0.3em}
  \fuente{Congreso de la Unión, minuta de reforma a la LISR 2027; LIF 2027 (transitorio vigésimo primero), aprobada por el Senado el 27 de octubre de 2026. Pérdidas y regularización: minutas aprobadas por la Cámara de Diputados en octubre de 2026.}
\end{frame}

% ===== A2. Cuatro frentes de reforma (E9) =====
\begin{frame}{Apéndice · Cuatro frentes de reforma}
  \vspace{0.3em}
  \begin{columns}[T,onlytextwidth]
    \begin{column}{0.49\textwidth}
      \caja{Reforma fiscal clásica (federal)}{Bases gravables, gastos fiscales, IVA, ISR.}\par\vspace{0.5em}
      \caja{Haciendas locales}{Ingresos propios de estados y municipios (predial, tenencia moderna) y financiamiento de sus funciones.}
    \end{column}
    \begin{column}{0.49\textwidth}
      \caja{Seguridad social}{Pensiones y salud con financiamiento sostenible e inclusión de trabajadores informales.}\par\vspace{0.5em}
      \caja{Empresas del Estado}{Pemex y CFE: reglas claras para los apoyos fiscales y registro transparente de los pasivos contingentes que hoy quedan fuera del presupuesto.}
    \end{column}
  \end{columns}
  \vspace{1.0em}
  {\small Ninguno de los cuatro frentes cierra la cuenta solo.\par}
\end{frame}

% ===== A3. Detalle del techo de endeudamiento =====
\begin{frame}{Apéndice · Detalle del techo de endeudamiento, LIF 2027}
  \vspace{0.4em}
  {\small
  \begin{tabularx}{\textwidth}{@{}>{\raggedright\arraybackslash}p{3.1cm}>{\raggedright\arraybackslash}X>{\raggedright\arraybackslash}X@{}}
    \toprule
    & \textbf{Endeudamiento neto interno} & \textbf{Endeudamiento neto externo}\\
    \midrule
    Gobierno Federal (techo) & Hasta 1.7 billones de pesos & Hasta 13,500 mdd\\\addlinespace
    Pemex & Hasta 160,619.5 mdp & Hasta 5,259.4 mdd\\\addlinespace
    CFE & Hasta 16,742 mdp & Hasta 725 mdd\\
    \bottomrule
  \end{tabularx}}
  \vspace{0.7em}
  \begin{itemize}\setlength{\itemsep}{0.3em}\small
    \item Gobierno Federal: los techos interno y externo son compensables entre sí. Suman alrededor de 1.94 billones de pesos, \dest{4.9\%} del PIB.
    \item El endeudamiento programado es de 1.57 billones de pesos, \dest{4.0\%} del PIB. La meta de RFSP es 3.9\% del PIB.
    \item El techo de Pemex está condicionado a su meta de balance financiero. Los pasivos contingentes no aparecen en el presupuesto ni en los techos de deuda.
  \end{itemize}
  \fuente{LIF 2027, arts. 1o, 2o y 4o; SHCP, CGPE 2027. \% del PIB con el PIB de SHCP.}
\end{frame}

\end{document}
